% Standalone section for incorporation into the main construction report.
% This is an ordinary mathematical proof, not a claim of a new Rocq endpoint.
\section{Changing the axiom constructor by eliminating double negations}
\label{sec:full-dn}

This section proves that the modified source searches exactly the same
numeric theorems as the original K230/K225 source.  The source programs
are not asserted to have the same register values at every boundary:
the code of the set-theory axiom changes.  Instead, the two numeric
calculi have identical sets of derivable codes.  In particular they
derive the halting code simultaneously.

\subsection{The numeric calculus, including its false code}

Use the original dual-Cantor pairing function $P$ and code constructors
\[
 I(A,B)=4P(A,B),\qquad Q_xA=4P(x,A)+2,\qquad
 M(x,y)=4P(x,y)+1,\qquad E(x,y)=4P(x,y)+3.
\]
Write $F=1$, $N A=I(A,F)$,
$A\mathbin{\&}B=N I(A,N B)$, and
$X_x A=NQ_xN A$.  These are abbreviations for functions on
\emph{all} natural numbers; no well-formedness restriction is imposed
on their arguments.  We use the named schematic letters only to make
their syntactic nesting readable.

Let $L$ be the numeric proof system retaining the original truth axiom
$0$, the eight rows $\rho_1,\ldots,\rho_8$, modus ponens and
generalization, but having no bundle axiom.  The following properties
of $L$ do not use a bundle instance.

\begin{enumerate}
\item The retained \L ukasiewicz row $\rho_8$ derives the implicational
axioms K, S and Peirce.  The explicit finite derivations are retained
in \texttt{coq/logic/CD230.v}.
\item $L\vdash I(F,A)$ for every $A$.  Indeed, K gives
$F\to X_0Q_0A$, because $X_0Q_0A$ has the form $B\to F$;
the retained row $\rho_6$ is $X_0Q_0A\to A$.  Transitivity gives
$F\to A$.  This is exactly the argument for
\texttt{thm\_F} in \texttt{Comp\_Prop230.v}.
\item Consequently $L$ proves both $A\to NN A$ and $NN A\to A$.
For the second direction, $NN A=(A\to F)\to F$ implies
$(A\to F)\to A$ by the preceding explosion theorem, and Peirce
with arguments $A,F$ finishes the proof.  This argument is uniform
even when $A=0$ or its arithmetic code has a special decoding case.
\item Mutual $L$-derivability is a congruence for $I$, $N$ and every
$Q_x$.  The first two assertions are propositional consequences of
K and S.  For the quantifier assertion, generalize a proof of
$A\to B$ and apply the unchanged distribution row
\[
 \rho_2(x,A,B):\quad Q_x I(A,B)\to I(Q_xA,Q_xB).
\]
Repeat for the reverse implication.  There is no free-variable side
condition here, since these are proofs of theorems and the numeric
generalization rule is unrestricted.
\end{enumerate}

It follows that deletion of any syntactically adjacent pair $NN$,
including below quantifiers, gives a mutually derivable numeric code.
This direct numeric argument is preferable to reasoning only about a
decoded formula: $F$ is literally the membership code $1$ in the
machine's arithmetic, and the theorem that it serves as falsity in
this proof system has just been supplied without using set theory.

\subsection{The two local constructor replacements}

The first replacement is
\begin{align*}
 X_x(A\mathbin{\&}B)
 &=NQ_x\bigl(NN I(A,N B)\bigr),\\
 J_x(A,B)&:=NQ_x I(A,N B).
\end{align*}
The two codes are mutually derivable in $L$.  On the source's postfix
stack, this changes the closed sequence
\[
 \texttt{and; ex}\quad\longmapsto\quad
 \texttt{neg; im; all; neg}.
\]
Starting with a top value $B$ and the two next stack operands $A,x$,
both sequences finish with precisely one value above the same lower
stack tail: the old value is $X_x(A\mathbin{\&}B)$, while the new
value is $J_x(A,B)$.  Each constituent macro terminates normally,
preserves the parameter and proof-list registers, and restores the
two scratch registers to zero.  These are the contracts proved for
\texttt{cons119}, \texttt{wal180}, \texttt{wn218},
\texttt{wa178}, and \texttt{wex183} in
\texttt{coq/prog/ProgLib230.v}.

The second replacement is
\begin{align*}
 A\mathbin{\&}(B\mathbin{\&}C)
 &=N I\bigl(A,NN I(B,N C)\bigr),\\
 T(A,B,C)&:=N I\bigl(A,I(B,N C)\bigr),
\end{align*}
again a pair of mutually derivable codes.  Its source realization is
\[
 \texttt{and; and}\quad\longmapsto\quad
 \texttt{neg; im; im; neg}.
\]
Starting with top $C$ and operands $B,A$, both words consume exactly
those two operands.  They preserve the entire lower stack tail and
the same clean-register boundary.  The following copy macro, used by
the existing freshness gate, is retained in the same position.

Binary-tree reassociations of these words are permitted because all
the displayed macro boundaries are normal exits: no break, continue
or halt escapes from a macro.  Internal loops are retained.  A
reassociation is not justified merely by its leaf order if it moves
an escaping control operation across a branch boundary.

\subsection{The complete changed bundle}

The following definitions give its exact arithmetic expression.
Let
\begin{align*}
 S(a,b,c)&=Q_c I\bigl(M(c,a),M(c,b)\bigr),\\
 R'(b,c,m)&=Q_b I\bigl(m,J_b(m,M(b,c))\bigr),\\
 H'(a,b,c,m)&=Q_b I\bigl(M(b,c),
   J_a(m,Q_a I(S(a,b,c),R'(b,c,m)))\bigr),\\
 K'(a,b,c,m)&=J_c\bigl(I(m,M(b,c)),H'(a,b,c,m)\bigr),\\
 V_M(a,b,c)&=I\bigl(M(a,b),Q_cM(a,b)\bigr),\\
 V_E(a,b,c)&=I\bigl(E(a,b),Q_cE(a,b)\bigr),\\
 B'(a,b,c)&=T\bigl(K'(a,b,c,c),V_M(a,b,c),V_E(a,b,c)\bigr).
\end{align*}
The original $K$ is obtained by replacing every $J_x(U,V)$ by
$X_x(U\mathbin{\&}V)$, and the original bundle is
$B(a,b,c)=K(a,b,c,c)\mathbin{\&}(V_M(a,b,c)\mathbin{\&}V_E(a,b,c))$.
Successively applying the first local equivalence to the witness
request, the local support block, and the outer closure block,
then applying the second equivalence to the final bundle,
proves, for all natural $a,b,c$,
\[
 L\vdash I(B(a,b,c),B'(a,b,c)),\qquad
 L\vdash I(B'(a,b,c),B(a,b,c)).                 \tag{DN}
\]
No binder moves, no variable is renamed, and no quantifier side
condition is introduced.  In particular, the reused $a$ and $b$
labels and their shadowing have exactly the original scopes.

The gate remains
\[
 G'(a,b,c)=
 \begin{cases}
  0,&c=a\text{ or }c=b,\\
  B'(a,b,c),&\text{otherwise}.
 \end{cases}
\]
On disabled inputs the old and new gates are literally equal; on
enabled inputs (DN) gives mutual derivability.  The new bundle is
an outer $N$ constructor, hence a multiple of four and therefore
different from $F=1$.  Thus $G'(a,b,c)\ne1$ on every input, including
the disabled cases.

\subsection{Proof translations and exhaustive enumeration}

Let $\mathcal C$ be the old numeric calculus and $\mathcal C'$ the
calculus that changes only its enabled bundle axiom from $B$ to $B'$.
Both contain $L$.  Induction on a finite proof in $\mathcal C$
translates it to a proof of the \emph{same numeric conclusion} in
$\mathcal C'$.  For a bundle-axiom line, use the corresponding $B'$
axiom and the second implication in (DN), followed by modus ponens.
Every other axiom and rule is copied unchanged.  The reverse
translation uses the first implication in (DN).  These are effective
finite syntactic proof transformations, formalizable in ZFC; in fact
their construction uses only elementary reasoning about finite
derivations.  Consequently
\[
 \forall n\quad
 \bigl(\mathcal C\vdash n\ \Longleftrightarrow\
       \mathcal C'\vdash n\bigr).             \tag{PT}
\]

The changed AX routine preserves the old proof-list reads, parameters,
generalized-cache formula, and the eight logical axioms pushed on the
stack.  Only its final top formula becomes $G'$.  Thus its boundary
contract is exactly the old \texttt{ax1} contract with \texttt{gate}
replaced by $G'$ and with \texttt{ax\_stack} unchanged.  The finite
postfix checks above establish that no extra truth cell is inserted
or deleted.  In particular this construction does not incorporate
the separate proposed MP-stack simplification.

The existing proof-enumerator soundness invariant transfers verbatim:
the current formula, cached formula, and all stack entries are
$\mathcal C'$-theorems.  An enabled AX produces a $\mathcal C'$ axiom,
a disabled AX produces its truth axiom $0$, and all other commands
retain their definitions.  Hence halting implies
$\mathcal C'\vdash1$.

For completeness, the existing proof-to-command compiler also
transfers.  Its bundle block selects $G'$ with the same parameters;
its other blocks depend only on the preserved stack positions,
logical rows, MP, DROP, NIP, and the facts that a disabled gate is
$0$ and every gate is different from $1$.  Both facts were proved
above.  Candidate decoding and the strictly draining proof-list
argument are unchanged.  The original bootstrap seed is
\[
 (\textit{prooflist},\textit{nextproof},\textit{top},
   \textit{stack},\textit{param1},\textit{cache})
       =(0,3,0,0,3,0),
\]
not the all-zero seed.  The original \texttt{START} command still
has selector at least $3$, so the constructed halting command list
has candidate number at least $3$ and is eventually examined from
this seed.  The retained seed-completeness argument therefore gives
\[
 \text{modified source halts from the inherited seed}
 \quad\Longleftrightarrow\quad \mathcal C'\vdash1.
\]
Combining this with (PT), the original K230 calibration, and a
separate universal binary-runtime certificate for the changed
transition table yields the desired equivalence with
$\neg\operatorname{Con}(\mathrm{ZFC})$.  The logical step here
is an ordinary proof formalizable in ZFC.  It is not advertised as
a newly compiled, axiom-free Rocq end-to-end theorem.

\subsection{Finite artifacts and the precise scope of checking}

The three old \texttt{and;ex} chunks have source prefixes
\texttt{1110100100111010111},
\texttt{11101001001110110101}, and
\texttt{111010010011101101101}.  The final
\texttt{and;(and;copy)} chunk has prefix
\texttt{111010010011101111011}.  These four changes remove eight
uses of the 47-leaf negation constructor, a total of 376 primitive
instruction occurrences: the source has 4745 leaves instead of 5121.
This is a source-size statement, not a Turing-state bound.

The companion \texttt{verify\_full\_dn\_semantics.py} checks the
exact pinned old source; the six reused macro bodies against their
K230 structured counterparts; the four replacement locations; the
new instruction word and all source guard successors; both symbolic
stack contracts; and equality of the old and new complete bundle
expressions after recursive double-negation removal.  Macro-body
comparison uses a distinguished normal-return marker and contracts
only finite runs of commuting increments and jumps, with every
decrement and branch retained.  Its checks are finite syntactic
obligations supporting the mathematical proof above.  They do not
replace the separate unbounded simulation check from the literal
binary Turing table to this source.
